\documentclass[10pt,a4paper]{amsart}
\usepackage[margin=19mm]{geometry}
\usepackage{amsmath,amssymb,mathtools}
\usepackage[scr=rsfs,cal=euler]{mathalfa}
\usepackage{xfrac}
\usepackage[unicode,hidelinks,bookmarks=false]{hyperref}
\newcommand{\ie}{i.e.\ }
\newcommand{\cf}{cf.\ }
\newcommand{\resp}{respectively\ }
\newcommand{\Subsection}{\subsection}
\providecommand{\nobreakdash}{\nobreak\hbox{-}}
\setlength{\emergencystretch}{2em}
\multlinegap0pt
\usepackage{bm}
\usepackage{bbm}
\usepackage{dsfont}
\usepackage{xspace}
\usepackage{cancel}
\usepackage{mathtools}
\usepackage{amsthm}
\makeatletter
\@ifundefined{thm}{\newtheorem{thm}{Thm}}{}
\@ifundefined{prop}{\newtheorem{prop}[thm]{Proposition}}{}
\makeatother
\newcommand{\R}{\mathbb{R}}                     % Reals
\title[The Liv\v{s}ic equation on differential forms over Anosov fl]{The Liv\v{s}ic equation on differential forms over Anosov flows and applications}
\author{Slobodan N. Simi\'c}
\date{Source: arXiv:2511.05678v1 (CC BY 4.0)}
\begin{document}
\maketitle

\noindent\emph{Source and licence.} The Liv\v{s}ic equation on differential forms over Anosov flows and applications. Version arXiv:2511.05678v1 (CC BY 4.0), distributed under the Creative Commons Attribution 4.0 International licence, as recorded in the arXiv licence metadata for this exact version and shown on the abstract page https://arxiv.org/abs/2511.05678v1. Full rights notice: licenses/arxiv-2511.05678-CC-BY-4.0.txt.\par\smallskip

\noindent\textit{Editorial scope.} Counted unit: the Prop whose source label is \texttt{prop:invariant} in the article (source0.tex lines 880--893, source label \texttt{prop:invariant}), delivered together with the article's own complete original proof of it (source0.tex lines 895--969, one \texttt{proof} environment of 75 lines). No mathematics is written, repaired or re-derived: statement and proof are the article's own text line for line. Required context retained uncounted: none. Every definition, display and label of those lines is retained unchanged. Omitted from this package and not redistributed: 1--879, 894--894, 970--1425. Nothing inside a retained excerpt is omitted or altered: every retained body is delivered exactly as the article's lines are written, so no omission receipt of any kind accompanies this package. Citations: the retained excerpts carry no citation command at all, so no bibliography entry of the article is transcribed and no reference is redistributed. Reference adaptation: none was needed and none was made; every reference inside the delivered unit resolves inside it, so no unresolved reference or citation is produced. Printed slips kept literal: no printed word, symbol, hypothesis or number of the counted statement or of its proof was altered. Layout and class: the article's own class is \texttt{scrartcl} with packages including amsmath, amsfonts, amssymb, amsthm, amscd, tikz, tikz-cd, hyperref, latexsym, mathrsfs, psfrag, epsfig, xy, wrapfig; the wrapper is the lane's pinned \texttt{amsart} editorial wrapper at 10pt on A4, which restates the theorem-like environments the retained text needs and adds the article's own macro definitions that the retained text needs (\texttt{\textbackslash R}), with extra wrapper packages (bm, bbm, dsfont, xspace, cancel, mathtools). The delivered numbering is the unit's own. Assets: no retained span contains a figure environment, a graphics-inclusion command, a PDF-page inclusion, a table or a TikZ picture; no image file is copied into this package, the package declares no asset dependency and no image file is redistributed with it. Licence: version arXiv:2511.05678v1 of this article is distributed under the Creative Commons Attribution 4.0 International licence, as recorded in the arXiv licence metadata for this exact version and shown on the abstract page https://arxiv.org/abs/2511.05678v1. A full rights notice is delivered at licenses/arxiv-2511.05678-CC-BY-4.0.txt. Characters: every retained excerpt is pure ASCII, so the delivered engine, XeLaTeX with its default Latin Modern text font, reports no missing character.
\begin{prop}   \label{prop:invariant}
  Let $\Phi$ be a smooth Anosov flow with infinitesimal generator
  $X$, preserving a smooth volume form $\Omega$. Then:

  \begin{enumerate}
  \item[(a)] $\text{\normalfont Inv}^0(M,X)$ consists of constant functions.
  \item[(b)] $\text{\normalfont Inv}^1(M,X) = \R \alpha$.
  \item[(c)] If the flow is asymmetric and $2 \leq k \leq n-2$, then
    $\text{\normalfont Inv}^k(M,X) = \{ \mathbf{0} \}$.
  \item[(d)] $\text{\normalfont Inv}^{n-1}(M,X) = \R \: i_X \Omega$.
  \item[(e)] $\text{\normalfont Inv}^n(M,X) = \R \Omega$.
  \end{enumerate}
  
\end{prop}

\begin{proof}
  %
  (a) and (e) are clear. To prove (b), assume $L_X \omega = 0$, for
  some $\omega \in C^1_X \Lambda^1(M)$.  Then
  $f_t^\ast \omega = \omega$, for all $t$, which clearly implies that
  $\omega(v) = 0$, for all $v \in E^{ss} \oplus E^{uu}$. Thus
  $\omega = \psi \alpha$, for some continuous $\psi : M \to \R$. Since
    %
    \begin{displaymath}
      0 = L_X \omega = (X\psi) \alpha = \psi L_X \alpha = (X \psi) \alpha
    \end{displaymath}
    %
    it follows that $\psi$ is flow invariant. Since the flow is
    ergodic (being volume preserving), $\psi$ is constant a.e., hence
    constant by continuity.

    (c) Assume the flow is asymmetric, $2 \leq k \leq n-2$, and
    $L_X \eta = 0$, for a continuous $k$-form $\eta$. We again have
    $f_t^\ast \eta = \eta$, for all $t$. Let $v_1, \ldots, v_k$ be
    arbitrary linearly independent vectors in the same tangent space
    of $M$. We claim that $\eta(v_1,\ldots,v_k) = 0$. Since
    $TM = E^{cs} \oplus E^{uu}$ and $\eta$ is multilinear, it is
    sufficient to show $\eta(v_1,\ldots,v_k) = 0$ in the following two
    cases:

  \begin{enumerate}
%
  \item[\textbf{Case 1}:] $\{ v_1, \ldots, v_k \} \subset E^{cs}$.
%
  \item[\textbf{Case 2}:]
    $\{ v_1, \ldots, v_k \} \subset E^{uu} \cup E^{cs}$ and at
    least one vector $v_j$ is in $E^{uu}$.
%
  \end{enumerate}

  \nopagebreak[4] In Case 1, by decomposing each $v_j$ into the sum
  $v_j = v_j^c + v_j^s \in E^c \oplus E^{ss}$, using the flow
  invariance of $\eta$, and the fact that $k \geq 2$, we obtain
%
    \begin{displaymath}
      \eta(v_1,\ldots,v_k) = \eta(f_{t \ast}(v_1), \ldots, f_{t
        \ast}(v_k)) \to 0,
    \end{displaymath}
%
    as $t \to +\infty$.

    In Case 2, the asymmetry of the flow implies 
%
    \begin{displaymath}
      \eta(v_1,\ldots,v_k) = \eta(f_{t \ast}(v_1), \ldots, f_{t
        \ast}(v_k)) \to 0,
    \end{displaymath}
%
    as $t \to -\infty$. Thus $\eta = 0$, as desired.

    To prove (d), assume $L_X \Theta = 0$, for some 
    $\Theta \in C^1_X \Lambda^{(n-1)}(M)$. Observe that since
    $L_X (i_X \Theta) = i_X L_X \Theta = 0$\footnote{Note that $i_X$
      and $L_X$ do commute on $C^1_X \Lambda^\ast(M)$.}, $i_X \Theta$
    is a continuous invariant $(n-2)$-form, hence zero by (c).

    Consider the continuous $n$-form $\alpha \wedge \Theta$. Since
    $L_X(\alpha \wedge \Theta) = L_X \alpha \wedge \Theta + \alpha
    \wedge L_X \Theta = 0$, $\alpha \wedge \Theta$ is invariant, hence
    $\alpha \wedge \Theta = c \: \Omega$, for some constant $c$. It
    follows that
    %
    \begin{displaymath}
      \Theta = i_X (\alpha \wedge \Theta) = i_X (c \: \Omega) = c \:
      i_X \Omega,
    \end{displaymath}
    %
    as desired.
    %
\end{proof}

\end{document}
